% ---
% Reemplaza líneas 1230-1274 de Docs/Technical_Report.tex
% Decisiones: N/A (sin cambio)
% ---

\subsubsection{El \textit{feature map} como mapeo no lineal al espacio de Hilbert}
\label{sec:feature_map}
% -----------------------------------------------------------

El bloque $U(\mathbf{x})$ define un mapeo no lineal:

\begin{equation}
  \phi: \mathbb{R}^k \rightarrow \mathcal{H}, \qquad
  \mathbf{x} \mapsto |\phi(\mathbf{x})\rangle =
  U(\mathbf{x})\,|0\rangle^{\otimes n}
  \label{eq:feature_map}
\end{equation}

\noindent donde $\mathcal{H} = \mathbb{C}^{2^n}$ es el
espacio de Hilbert del sistema de $n$ qubits, de dimensión
$2^n$. Como se argumentó en la sección del planteamiento
del problema, este espacio no es simplemente una versión
ampliada del espacio clásico de entrada: su geometría está
determinada por productos internos cuánticos
$\langle\phi(\mathbf{x})|\phi(\mathbf{x}')\rangle$, no por
distancias euclidianas. Con compuertas de entrelazamiento,
el mapeo $\phi$ hace depender el estado de productos de pares
de componentes de $\mathbf{x}$ además de cada componente por
separado \cite{Schuld2019QuantumML, havlicek2019supervised}; si
la geometría resultante aporta separabilidad frente a un kernel
clásico es una pregunta empírica, que es la que aborda este trabajo.
% Corregido el 2026-10-19 (auditoría A5-1): decía «correlaciones de orden superior que los métodos
% clásicos no capturan», en contradicción con §2.1.3 y §3.2.4 (interacciones de dos cuerpos como máximo).

Este mapeo define implícitamente un \textit{kernel cuántico}:

\begin{equation}
  k_Q(\mathbf{x}, \mathbf{x}') =
  \left|\langle\phi(\mathbf{x})|\phi(\mathbf{x}')\rangle
  \right|^2 =
  \left|\langle 0|^{\otimes n}\,U^\dagger(\mathbf{x}')\,
  U(\mathbf{x})\,|0\rangle^{\otimes n}\right|^2
  \label{eq:quantum_kernel}
\end{equation}

\noindent La elección del \textit{feature map} $U(\mathbf{x})$
determina completamente la geometría del espacio de
representación y, con ello, la capacidad del sistema para
separar las clases en ese espacio. En el presente trabajo
se comparan dos diseños de \textit{feature map}:
\textit{ZZFeatureMap} y \textit{PauliFeatureMap},
descritos a continuación.
